\chapter{The remaining opening when the opposite event is bounded}
\label{component:pd7bd7f0d7fb6}
\noindent\textbf{Component }\path{pd7bd7f0d7fb6}\par
{\small\noindent Source: \path{R13/baseline/baseline/baseline/baseline/baseline/advances/c2/c2_last_opening.tex}\par}
{\small\noindent SHA256: \path{d7bd7f0d7fb66e4e45a7fceaac4ff793b5d046b7874bc3879cd8bd20219c3c81}\par}
{\small\noindent Inventory use: CG1--CG10 complete gInf,runit opening, triangular m hard normal and literal n=0\par}
\medskip
\section{The remaining opening when the opposite event is bounded}
\label{pd7bd7f0d7fb6:sec:c2-last-opening}

\begin{theorem}[The unbounded-$m$ opening with bounded $n$]
\label{pd7bd7f0d7fb6:thm:c2-no-opening-bounds}
In Theorem~\ref{p81893ad08bed:thm:c2-three-openings}, replace the assumption that
$g=b-1+r+xh$ is a positive bounded unit by $g>0$ and the assumption that at least one of $r,g$ is a bounded
unit.  The all-tests tensor-SOS conclusion still holds.
\end{theorem}

The only new case is infinitesimal $g$; it is automatically bounded.
Our hypothesis now makes $r$ a unit.  Since $b_0+r_0=1$ and $b$ is
a positive unit, $r_0<1$.  Thus $n$ is bounded,
$0<b_0<1$, and $1-r$ is a unit.  The source formulas imply
\[
               mg=\frac{(1+k^2)(1-r)A_cS_c}{2\Delta},   \tag{CG1}
\]
which is a positive bounded unit.  Thus $m$ is unbounded.  The
two-opening bounds (CO1), positivity of $D/x,\alpha/x$, and the
consequence that $b-k^2r$ is a unit remain valid.

\paragraph{Use the original strong amplitude as the hard normal.}
Put $\eta=h/b$ and use the exact contact graph
\[
 Y_g=\eta uW,\qquad
 X_g=v(1+\eta u^2)W-(1+g)uZ.                            \tag{CG2}
\]
At contact this gives $X_g=\rho[b+xh-1-g]=-\rho r$.
The original $m$ amplitude satisfies exactly $M|_g=guZ$.
Normalize the full hard directions by
\[
 E_M=\sqrt m\{X-X_g-uv(Y-Y_g)\},\qquad
 E_Y=\sqrt{C+A}(Y-Y_g).                                 \tag{CG3}
\]
In these variables
\[
 M=guZ-E_M/\sqrt m,
 \quad
 X=X_g+E_M/\sqrt m+uvE_Y/\sqrt{C+A}.
\]
Consequently the unbounded packet has no $E_Y$ contribution at all.
This triangular change is essential: scaling $X$ alone would leave
the uncontrolled coefficient $\sqrt{m/(C+A)}uv$.
The entire scalar contribution of that packet is $4mg^2u^2Z^2$,
which is bounded by (CG1).  Its full $X$ derivative contributes
$-8mg\,uZ$, also bounded, before division by $\sqrt m$.

All other source weights entering the $X$ row are $n,D$ and the
unit signed weight, so the full $X$ derivative on (CG2) has bounded
coefficients on $V=vW,B=uZ$.  Hence the $E_M$ first normal has
zero-residue bounded coefficients on those roots.  The $E_Y$ first
normal is
\[
       \frac{(F_Y+uvF_X)|_g}{\sqrt{C+A}}.               \tag{CG4}
\]
The $m$ packet cancels identically in $F_Y+uvF_X$, before any graph
substitution.  The remaining large terms are exactly the $C,A$
terms controlled by (CO1); all other coefficients are bounded.
In particular the old (CO5) form on $P,B,E$ holds with bounded
coquotients and identically zero leading corrections.
The complete quadratic hard block is bounded and has residue
\[
                         \operatorname{diag}(4,4D_{v,0})  \tag{CG5}
\]
on the finite $u=0$ axis.  The $m$ packet gives precisely its first
diagonal entry, the $C,A$ packets give the second, and all mixed
entries from bounded packets contain $1/(\sqrt m\sqrt{C+A})$.

\paragraph{The remaining five Laurent terms already force strictness.}
The exact scalar divisions on $P,B,V,E$ are bounded for the same
reason: $C,A$ enter through (CO1), and $m$ only through $mg^2$.
Their leading weighted scalar is
\[
 q\,u^2W^2/x+b_su^2Z^2+c_sW^2-d_suWZ+exZ^2,
\]
where, for $z=v^2$ and $\nu=2n_0+1\geq1$,
\[
 \begin{gathered}
 q=Q(1+z),\quad b_s=2\nu(1+z)-2,\quad
 c_s=2z\nu(1+z),\quad d_s=4v\nu(1+z),\\
 e=E_2z+E_0,\qquad Q=4T_0/b_0,
 \quad E_2=4(D/x)_0>0,\quad E_0=4(\alpha/x)_0>0.
 \end{gathered}                                         \tag{CG6}
\]
Set $\mathcal D(s,z)=(qs+c_s)(b_ss+e)-d_s^2s/4$.
Its exact Laurent identity is
\begin{align*}
 \frac{\mathcal D(s,z)}{(1+z)sz}={}&
 2Q(\nu-1)s/z+2Q\nu s
 +2\nu E_2z/s+2\nu E_0/s\\
 &+QE_0/z+QE_2-4\nu.                                  \tag{CG7}
\end{align*}
Contact and stationarity give a zero at $(s,z)=(1,k_0^2)$.
After $z=k_0^2t$, write the five coefficients in order as
$a,b,c,d,e$.  Here $a\geq0$ and $b,c,d,e>0$.  Stationarity gives
$c=a+e$ and $b=d+e$.  The zero value fixes the constant, so exactly
\[
 \mathcal L=
 a(s/t+t/s-2)+d(s+1/s-2)+e(t/s+s+1/t-3).                 \tag{CG8}
\]
Every summand is nonnegative by AM--GM.  The last bracket alone
forces $s=t=1$ at equality: its three positive terms have product
one, and equality makes all three equal.  Its exponent vectors
$(-1,1),(1,0),(0,-1)$ span two dimensions; equivalently its Hessian
in $\log s,\log t$ at $(1,1)$ is
$\left(\begin{smallmatrix}2&-1\\-1&2\end{smallmatrix}\right)$.
Hence it already has a positive definite contact Hessian.
This proves global uniqueness and local nondegeneracy even when
$n_0=0$ and the first pair is absent.

For remote finite nonzero $v$ the same scalar cross-term split as
before gives a strict Gram.  The case $v_0=0,n_0>0$ is also strict
on $P,B,V,E$.  The full normals are restored by (CG3)--(CG5).
At matched finite-angle residues the whole scalar still has the
form $p^2Q_b/x+B_b$ with bounded $Q_b,B_b$.  The four monic
conductors $R,S,U,T$ of the prior theorem therefore give a bounded
Gram, and (CG8) proves its strict Schur inequality.

The full first normals have the required same conductor lifts.
For $E_M$, their coefficient is $(uA_1+WB_1)/\sqrt m$ with bounded
polynomial $A_1,B_1$.  Monic contact division retains the factor
$1/\sqrt m$ in every coquotient.  Exact contact makes the residual
coefficient $A_1+\kappa B_1$ vanish, giving the old four-root lift
with zero residue.  For $E_Y$, the cancellation in (CG4) leaves
precisely the old large-packet term $B_2$ plus bounded residual
terms; its exact contact and identically zero leading part give
the same $R B_2+\rho E_2'$ division, then the bounded factor
$1/\sqrt H$.  These are full normal identities, not just Hessian
checks.  Quadratic state degree and (CG5) finish the complete
matched-contact certificate.

\paragraph{The double-zero test when $n_0=0$.}
The graph (CG2) is replaced in this remote corner.  Put
\[
 H_c=4m+4n+2,\quad \xi=4m/H_c,\quad
 \zeta=(4m+2)/H_c,\quad \eta_c=C/(C+A).
\]
Use
\[
 Y_g=\eta_cuW,\quad
 X_g=\xi v(1+\eta_cu^2)W-\zeta uZ,
\]
and the triangular normals
\[
 E_X=\sqrt{H_c}\{X-X_g-uv(Y-Y_g)\},\qquad
 E_Y=\sqrt{C+A}(Y-Y_g).                                \tag{CG9}
\]
The entire graph amplitude is
\[
 M|_g=(1-\xi)v(1+\eta_cu^2)W+(\zeta-1)uZ.
\]
Both $m(1-\xi)$ and $m(\zeta-1)$ are bounded because $n$ is
bounded.  Thus the scalar and full first normals have bounded
coquotients after the same $x$ scaling as before.

Set $\epsilon=4\alpha+4n\zeta>0$.  It is infinitesimal, with
$n/\epsilon,x/\epsilon$ bounded.  The complete scalar $Z^2$
coefficient is exactly
\[
 f_{zz}=4\alpha+\epsilon u^2+4n\zeta^2u^4+4Dv^2
                            +u^2v^2 B_b,               \tag{CG10}
\]
where $B_b$ is bounded with residue $2$.  This identity follows
directly from the packets and the exact relation
$4n\zeta^2+(4m+2)(1-\zeta)^2=4n\zeta$.
The $W^2$ coefficients on $P=uW/\rho,V=vW$ have positive residues
$Q,2$.  The cross term is $uvWZ$ times a bounded function and
can be written as a zero-residue multiple of $PE$ using
$uvWZ=vPE$.  Thus the five scalar roots
\[
           P,\quad V,\quad E=\rho Z,\quad
           B_\epsilon=\sqrt\epsilon\,uZ,\quad T=uvZ
\]
have strict diagonal residue $(Q,2,E_0,1,2)$.

For the full $X$ derivative, at $v=0$ the $Z$ coefficient is exactly
$-8n\zeta u^3$.  After division by $\sqrt{H_c}$ its coefficient
on $B_\epsilon$ is infinitesimal because
$n/\sqrt\epsilon=(n/\epsilon)\sqrt\epsilon$ is infinitesimal.
The remaining $Z$ coefficient has the factor $uv^2$ and bounded
coquotient; its coefficient on $T$ has zero residue.  The $W$
coefficient has a factor $v$ and bounded coquotient, multiplied by
$1/\sqrt{H_c}$.  For $E_Y$ the exact $m$ cancellation in
$F_Y+uvF_X$ again applies; (CO6) gives
$B_2=8D_v(xA)\eta_c u^2$, and the remaining lift is
\[
 P(B_2+xe_w)+T(\rho u e_z)+E\,8v(D-\alpha),
\]
divided by $\sqrt H$, with bounded $e_w,e_z$.
Every coefficient has zero residue.  The hard residue of (CG9) is
$\operatorname{diag}(1,4)$, and all mixed hard entries vanish in
residue.  This proves the full-parent corner including literal $n=0$.

\paragraph{All other tests.}
For the remaining commuting axes use the earlier outer identities,
cap $C,A$ at $\sigma/x$ as before, and also cap $m$ at
$\min(m,1)$.  Retain the entire nonnegative excess $m$ packet.
The selected $m$ weight is a unit and all previous displayed
strict residues apply.  The noncommuting-residue theorem covers
the other tests.  Restore all positive factors and split complex
states as before.  This proves the theorem for arbitrary positive
$g$, including the new unbounded-$m$ regime.

The remaining assumptions in this source family are now only that
$k,b$ are units, $x$ is positive infinitesimal, all original contact
openings are positive, at least one of $r,g$ is a unit, $n\geq0$,
and the source is physical.  The joint degeneration $r,g\to0$,
where both $n,m$ can be unbounded, is not covered.  The
result does not yet treat unbounded $b$ or a degenerate residue of
the finite actual angle $k$.  Section~\ref{p7294aae40ee2:sec:c2-bounded-state-entry}
excludes infinitesimal $b$ under these entry hypotheses.
